\scanpage{169}
% archive sheet number: 168
Je dis que cette loi de composition fait de $E$ un \uline{schéma en groupes}. Il faut vérifier

(i) L'associativité, qui s'exprime par l'égalité de \uline{deux morphismes} $E\times E\times E\longrightarrow E$ au dessus de $G\times G\times G\longrightarrow G$, compatibles avec $\unsure{\gamma|\sigma}\times\gamma\times\gamma\longrightarrow\gamma$ ; cela résulte du fait qu'il y a \uline{un seul} tel morphisme, à translation près par un élément de $\operatorname{Hom}_S(G\times G\times G,\gamma)\simeq\gamma(S)$, donc un seul pour lequel l'unité entre en l'unité.

(ii) $e$ est l'unité, même argument, \deleted{\unsure{pour}} \unsure{prouver} que deux hom $E\longrightarrow E$ sont l'identité.

(iii) Existence d'une symétrie de $E$

Voir l'exposé du livre de Serre, page 183.

\noindent ----

\noindent\rule{.72\linewidth}{.3pt}

Application au cas où $\gamma$ sur $\deleted{E}S$ est \uline{étale}, \inserted{et $G$ à fibres géom. conn.}

\uline{Alors}
\[
\operatorname{Ext}^{1}_{S\text{-gr}}(G,\gamma)\unsure{\simeq|\cong}\operatorname{Ext}^{1}_{S_0\text{-gr}}(G_0,\gamma_0)
\]
où $S_0=S_{\mathrm{réd}}$, $G_0=G\times_S S_0$, $\gamma_0=\unsure{G|\gamma}\times_S S_0$. Ceci est plus gén. si $S_0\longrightarrow S$ est un homéom. universel quelconque. \deleted{Ceci permet l'application}

% slanted
\margin{N.B. Ce résultat est valable de façon indépendante de la prop. plus haut, si on suppose $G$ à fibres connexes --------}
